\section{Conclusions}
\label{sec:conclusion}

TRACT counts white-tailed deer from thermal road transects, and decomposing its error moves
the problem rather than solving it: what a practitioner would try to buy with a better
detector is lost after detection, in the stage deciding whether a track is a new animal.
Four independent improvements to detection and association each left the count no better or
worse, and learned confirmers of three capacities lost to three hand-tuned parameters. The
reason is supervision rather than model class: with roughly two hundred positive tracks, and
duplicates of animals already counted sitting \emph{between} the positives and the false
candidates on every feature we measured, what separates a first sighting from a re-sighting
is precisely what a confirmer cannot learn at this scale. Two boundaries follow for anyone
adopting the result. The pipeline does not transfer unaltered, since a detector trained
without a site scores it low enough that an absolute confidence threshold can reject every
track (\cref{sec:loso}). And the under-count is an operating-point choice rather than a
limit: at the threshold where signed bias crosses zero the same pipeline recovers 81 of the
83 held-out animals in aggregate (\cref{sec:operating}). We release the dataset,
per-individual annotations, evaluation harness and decomposition tooling; the gap between
what is reachable and what is counted is large, precisely localized and entirely open.
